\documentclass[%
biblatex,
%marginversioninfo
]{ThermodBook}


\NewTColorBox{hsolution}{+O{}}{%
  breakable, enhanced,sharpish corners, colframe=blue!50,colback=white,coltitle=blue!50!black,
  fonttitle=\itshape\bfseries, parbox = false,
  boxrule=0.3mm,top=0mm,bottom=0mm,
  code={\if@twoside\tcbset{check odd page, toggle enlargement, toggle left and right}\fi},
  left=1mm,
  right=\lsize,
  grow to right by=\lsize,
  underlay={%
    \begin{tcbclipinterior}
      \draw[help lines,step=5mm,blue!20!white,shift={(interior.north west)}]
      (interior.south west) grid (interior.north east);
     \if@twoside%
	     \ifoddpage%
	        \draw[red!75!black,line width=1pt]
	        ([xshift=-\lsize]frame.north west)
	        --([xshift=-\lsize]frame.south west);%
	      \else%
	        \draw[red!75!black,line width=1pt]
	        ([xshift=\lsize]frame.north east)
	        --([xshift=\lsize]frame.south east);%
	      \fi%
    \else
    \draw[red!75!black,line width=1pt]
      	 ([xshift=-\lsize]frame.north west)
     	  --([xshift=-\lsize]frame.south west);%
      \fi
    \end{tcbclipinterior}},
  title={\solut},
  label={},
  attach title to upper=\quad,
  after upper={\par\hfill%
      {}},
  lowerbox=ignored,
  #1
  }


\begin{document}

%=========================================================
\begin{problem}%
Розібрати, як буде поводитися при різних температурах від $0$ до $10^\circ$C термометр, заповнений
водою. Для яких температур показання цього термометра будуть однаковими? Для об'єму води в залежності
від температури можна прийняти формулу:

\[
V = 1 - 0.00006105\,t + 0.000007733\,t^2
\]

де $V$ -- об'єм при температурі $t$. Об'єм при $0^\circ$C прийнятий за одиницю.
\end{problem}


\section*{Розв'язок}

1. \textbf{Аналіз поведінки води:}

Вода має аномальні теплові властивості - її об'єм зменшується при нагріванні від
$0^{\circ}$C до $4^{\circ}$C, і збільшується при подальшому нагріванні.

2. \textbf{Знайдемо температуру мінімального об'єму:}

Знаходимо похідну об'єму за температурою:
\[
\frac{dV}{dt} = -0.00006105 + 2 \times 0.000007733\,t
\]

Прирівнюємо до нуля для знаходження екстремуму:
\[
-0.00006105 + 0.000015466\,t = 0
\]
\[
t = \frac{0.00006105}{0.000015466} \approx 3,95^{\circ}\text{C}
\]

3. \textbf{Однакові показання термометра:}

Термометр покаже однакові значення для температур, при яких об'єм води однаковий. Знайдемо $t_1$ і
$t_2$, для яких $V(t_1) = V(t_2)$.

Розв'яжемо рівняння:
\[
1 - 0.00006105\,t_1 + 0.000007733\,t_1^2 = 1 - 0.00006105\,t_2 + 0.000007733\,t_2^2
\]
\[
0.000007733(t_1^2 - t_2^2) - 0.00006105(t_1 - t_2) = 0
\]

Оскільки $t_1 \neq t_2$, ділимо на $(t_1 - t_2)$:
\[
0.000007733(t_1 + t_2) - 0.00006105 = 0
\]
\[
t_1 + t_2 = \frac{0.00006105}{0.000007733} \approx 7.9
\]

Таким чином, показання термометра будуть однаковими для будь-яких двох температур, сума яких дорівнює
приблизно $7.9^{\circ}$C, наприклад:
\begin{itemize}
\item $2^{\circ}C$ і $5.9^{\circ}$C
\item $3^{\circ}C$ і $4.9^{\circ}$C
\end{itemize}

4. \textbf{Висновок:}

Термометр, заповнений водою, буде:
\begin{itemize}
\item При $t < 3.95\textsuperscript{\circ}C$ --- зменшувати свій об'єм при нагріванні
\item При $t > 3.95\textsuperscript{\circ}C$ --- збільшувати свій об'єм при нагріванні
\item Показувати однакові значення для температур, сума яких дорівнює $7.9^{\circ}$C
\end{itemize}


\end{document}

